\section{Coverage, reproducibility, and remaining work}
\label{sec:coverage}

\subsection{What the packet establishes}

There are 71 numbered results in the geometric and analytic chapters:
30 in affine and metric geometry, 24 in conics and quadrics, nine in
coordinate systems, and eight in polynomial tangency and arc length.
Seven further foundational results specify the arithmetic and representation
discipline; one of those is an external computability obstruction, explicitly
not an internal consistency theorem. The total of 78 counts named theorem,
lemma, proposition, and corollary blocks. It does not count definitions,
examples, individual subclaims, or Lean certificates a second time.

\begin{center}
\begin{tabular}{p{.23\linewidth}p{.65\linewidth}}
\toprule Source area & Coverage and boundary \\\midrule
Cartesian geometry & Exact incidence, determinants, affine maps, distances,
projections, metric inequalities, and finite triangle/circle constructions.\\[4pt]
Conics & Complete affine classification for rational coefficients, explicit
degeneracies, focal equivalences, parametrizations, tangents, polars, optical
laws, and asymptote bounds.\\[4pt]
Quadrics & Rational elimination to diagonal/null forms, affine type table,
stereographic sphere coordinates, two ruled examples, ellipsoid support.\\[4pt]
Coordinates & Finite-series radian functions and forward polar/cylindrical/
spherical formulas; inverse direction coordinates with rational case tests
or supplied axis and gap data.\\[4pt]
Tangency & Polynomial remainder estimates; regular implicit curves and
surfaces constructed by bounded bisection; tangent errors and polynomial
arc-length approximation.\\\bottomrule
\end{tabular}
\end{center}

The Euler-line and nine-point-circle deductions, the local implicit theorem,
and the arc-length modulus are supplementary developments in these linked
topics. We do not pretend that each is stated verbatim in the parent article.
The machine-readable inventory marks them accordingly.

\subsection{Limits of the proved claims}

This packet does not claim a formalized proof of all elementary geometry,
general real-algebraic quantifier elimination, an algorithm to classify conics
with arbitrary real-name coefficients without promises, a global continuous
scalar angular inverse, or a theorem on every continuous curve. We do not prove
the general bound on the number of intersection points of two conics here;
that bound requires exclusion of common components and a separate elimination
argument. The line--conic, line--circle, and circle--circle results are complete
in their stated domains. The factorized equations of singular crossings do not
assert an unrestricted disjunction selecting one component for an arbitrary
represented-real input.

Every geometric assertion has a fully developed mathematical proof in its
stated E/A/R interpretation. The new complete object-language derivations are
still a formalization task. In particular, the 24 explicit polynomial Lean
candidates have not been compiled, and the rational/name bridge records have
no asserted instances. The absence of a proof placeholder in source code is
not itself a verification result.

\subsection{Executable evidence}

The exact-rational conic classifier returns an affine change, transformed
coefficients, and a type. Its independent checker reconstructs all six
coefficients by evaluating at $(0,0),(1,0),(-1,0),(0,1),(0,-1),(1,1)$.
These values determine a polynomial of total degree at most two, so this is
an exact polynomial-identity test for the specified change of coordinates.
The finite regression run covers all $5^6=15{,}625$ integer coefficient
sextuples in $[-2,2]^6$, all 12 returned classes, and rejection of a deliberately
altered certificate. This verifies those finite instances and the checker
behavior; it is not a proof of the classifier for all inputs.

Additional exact tests check the geometric formulas, square-root brackets,
trigonometric tails, Taylor remainders, implicit-root intervals, and chord
refinement bounds. An independent SymPy pass reads the arithmetic statements
in the candidate Lean file and checks 28 scalar identities across its 24
endpoints, plus 19 additional scalar identities in 13 formula families.
It checks the polynomial equations, not Lean syntax or proof terms; it cannot
establish the positivity guards or the constructive scope on its own.
The exact counts, source hashes, and software versions are recorded in
\path{artifacts/}. Floating-point plots are illustrative only and do not
decide any geometric predicate.

\subsection{Rebuilding the packet}

From the extracted packet directory, use
\begin{lstlisting}
python3 code/conic_classifier.py
python3 code/check_affine_metric.py
python3 code/check_analytic.py
python3 code/check_symbolic.py
python3 code/build_inventory.py
bash build.sh
\end{lstlisting}
The first three numerical programs use only Python's standard library.
The symbolic pass requires SymPy; rebuilding the figures requires Matplotlib
and NumPy. The report uses a standard TeX Live installation with
\texttt{latexmk} and \texttt{pdflatex}. The saved figures are already in the
archive, so regenerating them is optional. Lean checking has its separate
script and pinned dependency requirements in Section~\ref{sec:formalization}.
The archive includes source, validation receipts, and checksums for its files.

\subsection{The complete named-result ledger}

The following table is generated from the report's actual theorem blocks.
It is an index to the proofs, not a replacement for their hypotheses or
constructive interpretation. The topic names identify the source neighborhood;
the JSON inventory records proof status and supplementary results separately.

\small
\begin{longtable}{p{.08\linewidth}p{.55\linewidth}p{.28\linewidth}}
\toprule Result & Statement group & Source topic \\ \midrule
\endfirsthead
\toprule Result & Statement group & Source topic \\ \midrule
\endhead
\bottomrule\endfoot
\ref{found:rat} & Coded ordered-rational arithmetic & ProveIt iPRA framework \\
\ref{found:loops} & Finite loops and bounded formulas & ProveIt iPRA framework \\
\ref{found:sqrt} & Dyadic square-root bracket & ProveIt iPRA framework \\
\ref{found:holder} & Square-root order and a uniform modulus & ProveIt iPRA framework \\
\ref{found:operations} & Arithmetic on names with explicit precision & ProveIt iPRA framework \\
\ref{found:transfer} & Finite algebra and quantified name identities & ProveIt iPRA framework \\
\ref{found:no-zero-test} & No uniform zero test for all these names & ProveIt iPRA framework \\
\ref{geo:line} & Equivalent equations of a line, and the line through two points & Line (geometry) \\
\ref{geo:slope} & Slope, intercepts, parallelism, and perpendicularity & Slope; Line (geometry) \\
\ref{geo:two-lines} & Complete intersection classification of two rational lines & Line (geometry); Intersection (geometry) \\
\ref{geo:area} & Determinant, collinearity, and triangle area & Cartesian coordinate system \\
\ref{geo:section} & Section formula, midpoint, and segment inequalities & Line (geometry) \\
\ref{geo:centroid} & Barycentric coordinates and concurrence of medians & Affine coordinates \\
\ref{geo:affine} & Affine changes of coordinates & Cartesian coordinate system \\
\ref{geo:lagrange} & Lagrange's sum-of-squares identity & Dot product \\
\ref{geo:cauchy} & Cauchy--Schwarz and its equality condition & Dot product \\
\ref{geo:pythagoras} & Pythagoras, parallelogram identity, and polarization & Euclidean space; Dot product \\
\ref{geo:metric} & The Euclidean distance is a metric & Euclidean space \\
\ref{geo:cosine} & Cosine law and an algebraic encoding of planar angles & Angle; Dot product \\
\ref{geo:projection} & Projection onto a line and its minimum-distance property & Line (geometry); Dot product \\
\ref{geo:hyperplane} & Projection onto a line or plane in normal form & Plane (geometry); Dot product \\
\ref{geo:reflection} & Reflection in a line or plane & Cartesian coordinate system \\
\ref{geo:isometries} & Orthogonal transformations and similarities & Cartesian coordinate system \\
\ref{geo:cross} & Cross product, triple product, and volume & Cross product \\
\ref{geo:plane} & Equations of a plane and a plane through three points & Plane (geometry) \\
\ref{geo:line-plane} & Line--plane intersection, including all degenerate cases & Plane (geometry); Line (geometry) \\
\ref{geo:two-planes} & Intersection of two planes & Plane (geometry) \\
\ref{geo:skew} & Common perpendicular and distance between two spatial lines & Line (geometry); Cross product \\
\ref{geo:circle} & Circle equation and completion of the square & Circle \\
\ref{geo:line-circle} & Line--circle intersections & Circle; Intersection (geometry) \\
\ref{geo:circle-circle} & Radical axis and complete circle--circle intersection test & Circle; Intersection (geometry) \\
\ref{geo:circumcenter} & Circumcenter of a noncollinear rational triangle & Circle \\
\ref{geo:orthocenter} & Orthocenter and the Euler-line relation & Circle \\
\ref{geo:ninepoint} & The nine-point circle & Circle \\
\ref{geo:power} & Power of a point, secants, and equal tangents & Circle \\
\ref{geo:shoelace} & Shoelace formula and origin-independent polygon area & Cartesian coordinate system \\
\ref{geo:sphere-plane} & Intersection of a sphere with a plane & Quadric surface; Plane (geometry) \\
\ref{con:covariance} & Scaling and coordinate covariance & Conic section \\
\ref{con:invariants} & Discriminant and projective singularity & Conic section \\
\ref{con:completion} & Rational completion of squares & Conic section \\
\ref{con:classification} & Exhaustive real-locus classification for rational conics & Conic section \\
\ref{con:axes} & Constructive principal axes, circles, and rectangular hyperbolas & Conic section \\
\ref{con:roots} & Complete quadratic solver with degeneracies & Quadratic equation \\
\ref{con:line-intersection} & Line--conic intersection & Conic section \\
\ref{con:rational-point} & Parametrization from one rational point & Conic section \\
\ref{con:standard-param} & Standard rational parametrizations & Conic section \\
\ref{con:ellipse-foci} & Ellipse: focal and Cartesian definitions agree & Ellipse \\
\ref{con:hyperbola-foci} & Hyperbola: focal and Cartesian definitions agree & Hyperbola \\
\ref{con:parabola-foci} & Parabola: focus and directrix & Parabola \\
\ref{con:directrix} & Eccentricity and a guarded polar equation & Conic section; Polar coordinate system \\
\ref{con:tangent} & Quadratic tangents and second-order contact & Conic section \\
\ref{con:polar} & Polar reciprocity and poles & Conic section \\
\ref{con:ellipse-reflection} & Ellipse reflection law & Ellipse \\
\ref{con:parabola-reflection} & Parabola reflection law & Parabola \\
\ref{con:asymptotes} & Hyperbola asymptotes with a modulus & Hyperbola \\
\ref{con:diameters} & Parallel chord midpoints and conjugate diameters & Conic section \\
\ref{con:quadric-reduction} & Finite rational reduction of every quadric & Quadric surface \\
\ref{con:quadric-types} & Exhaustive affine types of rational quadrics & Quadric surface \\
\ref{con:sphere-param} & Stereographic rational parametrization of the sphere & Quadric surface; Parametric equation \\
\ref{con:rulings} & Two rulings of a hyperbolic paraboloid and a hyperboloid & Quadric surface \\
\ref{con:ellipsoid-support} & Ellipsoid support plane and unique maximizing point & Quadric surface; Dot product \\
\ref{coord:tail} & A primitive-recursive trigonometric tail bound & Polar coordinate system \\
\ref{coord:addition} & Addition laws and the unit-circle identity & Polar coordinate system; Rotation of axes \\
\ref{coord:derivatives} & Quantitative trigonometric derivatives & Polar coordinate system; Tangent \\
\ref{coord:polar} & Polar direction coordinates and reconstruction & Polar coordinate system \\
\ref{coord:cylindrical} & Cylindrical coordinates & Cylindrical coordinate system \\
\ref{coord:spherical} & Spherical coordinates with explicit axis data & Spherical coordinate system \\
\ref{coord:gram-metric} & The metric in an oblique coordinate frame & Cartesian coordinate system; Dot product \\
\ref{coord:locus-transform} & Transforming a locus uses the inverse substitution & Analytic geometry; Rotation of axes \\
\ref{coord:jacobians} & Coordinate Jacobian identities & Polar coordinate system; Cylindrical coordinate system; Spherical coordinate system \\
\ref{tan:taylor} & A polynomial Taylor estimate without a limit theorem & Tangent; Parametric equation \\
\ref{tan:graphs} & Tangent lines and planes to polynomial graphs & Tangent; Parametric equation \\
\ref{tan:parametric} & Polynomial parametrized curves and tangent transport & Tangent; Parametric equation \\
\ref{tan:divided} & Explicit divided-difference bounds & Tangent; Parametric equation \\
\ref{tan:implicit} & Quantitative polynomial implicit function & Tangent; Parametric equation \\
\ref{tan:implicit-tangent} & The tangent of the implicit branch, with its error & Tangent; Parametric equation \\
\ref{tan:singular} & Regularity is an essential hypothesis & Tangent; Parametric equation \\
\ref{tan:arclength} & Constructive arc length of a polynomial curve & Tangent; Parametric equation \\
\end{longtable}\normalsize

